\documentclass[a4paper]{article}
% Español (México) con XeLaTeX: fontspec para fuentes Unicode, polyglossia
% para la división silábica y los nombres del documento en la variante mexicana.
\usepackage{fontspec}
\usepackage{polyglossia}
\setdefaultlanguage[variant=mexican]{spanish}
% Los nombres chinos van como «pinyin (original)»: xeCJK da a los caracteres
% Han su propia fuente; sin ella XeLaTeX los omite con un aviso.
% FandolSong no cubre algunos caracteres de nombres (U+73FA); WenQuanYi Zen Hei sí.
\usepackage[AutoFallBack=true]{xeCJK}
\setCJKmainfont{FandolSong-Regular.otf}
\setCJKfallbackfamilyfont{\CJKrmdefault}{WenQuanYi Zen Hei}
% polyglossia no traduce los nombres de listings.
\AtBeginDocument{\providecommand{\lstlistingname}{}\renewcommand{\lstlistingname}{Listado}%
  \providecommand{\lstlistlistingname}{}\renewcommand{\lstlistlistingname}{Índice de listados}}
\usepackage{amsmath, amssymb}
\usepackage{graphicx}
\usepackage[margin=2.35cm]{geometry}
\usepackage[most]{tcolorbox}
\usepackage{booktabs}
\usepackage{tabularx}
\usepackage{longtable}
\usepackage{array}
\usepackage{float}
\usepackage{xcolor}
\usepackage{hyperref}
\usepackage{fancyhdr}
\setCJKmainfont[AutoFakeBold=2.4]{AR PL UMing CN}
\setCJKsansfont[AutoFakeBold=2.4]{Droid Sans Fallback}
\setCJKmonofont[AutoFakeBold=2.4]{Droid Sans Fallback}
\hypersetup{colorlinks=true, linkcolor=blue!55!black, urlcolor=blue!60!black}
\graphicspath{{./}{figures/}}
\setlength{\parskip}{0.55em}
\setlength{\parindent}{2em}
\setlength{\emergencystretch}{3em}
\renewcommand{\arraystretch}{1.18}
\newtcolorbox{knowledgebox}[1]{enhanced, colback=blue!5!white, colframe=blue!70!black, colbacktitle=blue!70!black, coltitle=white, fonttitle=\bfseries, title={#1}, sharp corners, breakable}
\newtcolorbox{importantbox}[1]{enhanced, colback=yellow!10!white, colframe=yellow!75!black, colbacktitle=yellow!75!black, coltitle=black, fonttitle=\bfseries, title={#1}, sharp corners, breakable}
\newtcolorbox{warningbox}[1]{enhanced, colback=red!5!white, colframe=red!70!black, colbacktitle=red!70!black, coltitle=white, fonttitle=\bfseries, title={#1}, sharp corners, breakable}
\pagestyle{fancy}
\fancyhf{}
\fancyfoot[C]{\thepage}
\fancyfoot[R]{\footnotesize\color{gray}\href{https://github.com/hqhq1025/ai-course-notes}{AI Course Notes}}
\renewcommand{\headrulewidth}{0pt}
\renewcommand{\footrulewidth}{0pt}
\newcommand{\videourl}{https://www.youtube.com/watch?v=SG90aehV3vU}
\newcommand{\repourl}{https://github.com/hqhq1025/ai-course-notes}

\begin{document}
\begin{titlepage}
\centering
{\Large Notas de clase del informe trimestral sobre modelos grandes\par}
\vspace{1.0cm}
{\huge\bfseries AI War, las dos grandes alianzas y el tercer paradigma: Online Learning\par}
\vspace{0.5cm}
{\Large Zhang Xiaojun conversa con Guang Mi (广密)\par}
\vspace{0.35cm}
{\large Elaboradas a partir del video público de Zhang Xiaojun Podcast, la transcripción generada en GPU, la estructura de capítulos y diagramas conceptuales propios\par}
\vspace{0.3cm}
{\large 2025-12-24\par}
\vspace{0.85cm}
\includegraphics[width=0.88\textwidth,height=0.42\textheight,keepaspectratio]{cover.jpg}\par
\vfill
\begin{tcolorbox}[width=0.92\textwidth, colback=black!2!white, colframe=black!55, sharp corners]
\textbf{Título del video}: 127. Conversación de fin de año del informe trimestral sobre modelos grandes: Guang Mi pronostica una AI War, dos grandes alianzas y un tercer paradigma, Online Learning\par
\textbf{Canal del video}: Zhang Xiaojun Podcast\par
\textbf{Duración del video}: 01:18:03\par
\textbf{Enlace del video}: \href{\videourl}{\nolinkurl{\videourl}}\par
\textbf{Nota sobre los materiales}: YouTube no ofrece subtítulos disponibles; el texto se elaboró a partir de la transcripción con faster-whisper large-v3 en CUDA, los capítulos del video, la portada y figuras didácticas propias. El video muestra una toma fija de entrevista o podcast, por lo que el texto no inserta repetidamente fotogramas de las personas.\par
\textbf{Página del proyecto}: \href{\repourl}{\nolinkurl{\repourl}}\par
\end{tcolorbox}
\end{titlepage}

\tableofcontents
\newpage

\section{Introducción: por qué este episodio debe leerse como un material de clase trimestral}

Este episodio es la conversación de fin de año del «Informe trimestral de modelos grandes a nivel mundial». No es una entrevista biográfica común, sino una revisión de cara a 2026 de la industria de la IA, de los paradigmas de modelos y de las líneas principales de inversión. Zhang Xiaojun (张小珺) se encarga de llevar las preguntas al plano del negocio, las plataformas y el emprendimiento; Guang Mi (广密), por su parte, condensa un conjunto de observaciones trimestrales en unas cuantas palabras clave: AI War, liderazgo alternado, los modelos fundacionales como plataformas de comercio electrónico generalistas, las dos líneas principales de OpenAI y Google, los tres paradigmas de preentrenamiento/RL/Online Learning y «el modelo es el producto, los datos son el modelo».

Por ello, este material no sigue los subtítulos originales frase por frase, sino que se reorganiza según una lógica didáctica en seis líneas principales. La primera explica por qué «AI Bubble» no basta para explicar esta ronda de inversión: la IA se parece más a una carrera armamentista que nadie puede darse el lujo de perder. La segunda trata las cuentas de ingresos de OpenAI y la disputa entre Google y OpenAI por los puntos de entrada. La tercera aborda el liderazgo alternado de GPT, Claude y Gemini y la analogía entre las empresas de modelos fundacionales y las plataformas. La cuarta trata el tercer paradigma, Online Learning. La quinta abarca los datos, la Agentic Web, los Neo Labs y la Robotics. La sexta trata las narrativas de China y Estados Unidos, los emprendedores de origen chino y las variables de largo plazo para los próximos tres a cinco años.

\begin{importantbox}{Tesis central de este episodio}
El juicio central de Guang Mi no es que «no hay burbuja», sino que «el marco de la burbuja es demasiado estrecho». La inversión en IA está impulsada a la vez por la defensa comercial, la estrategia nacional, la oferta de capacidad de cómputo, la competencia por el talento y las expectativas sobre el siguiente paradigma. Se parece más a una AI War: nadie está seguro de que las cuentas de corto plazo vayan a cuadrar, pero ninguno de los actores clave se atreve a levantarse de la mesa.
\end{importantbox}

\begin{knowledgebox}{Nota sobre la estrategia visual}
Este video no tiene slides, pizarrón, imágenes de artículos ni demostraciones de productos. Conforme al flujo de trabajo para pódcasts, el cuerpo del texto usa la portada solo para identificar la fuente y no vuelve a insertar las tomas fijas de la entrevista; todas las figuras del texto son diagramas conceptuales de elaboración propia que muestran relaciones estructurales, y las explicaciones detalladas se ubican en el texto, las tablas y los recuadros de Lectura de la figura.
\end{knowledgebox}

\subsection{Resumen de la sección}

El valor de EP127 está en condensar los temas de modelos, datos, Agent, capacidad de cómputo y emprendimiento que aparecieron una y otra vez en episodios anteriores en un juicio de fin de año: en 2026 no hay que fijarse en un benchmark aislado, sino en si cambia el paradigma, cómo se redistribuyen los puntos de entrada y quién logra cerrar el ciclo entre datos y producto.
\section{De la AI Bubble a la AI War: por qué la «burbuja» no basta como explicación}

Esta sección aborda primero el marco general de todo el episodio. Al inicio del programa la pregunta es por la «AI Bubble», pero la respuesta de Guang Mi (广密) es que «se parece más a una AI War». No es un juego de palabras, sino un cambio en la unidad de análisis. Si la IA se ve solo como un proyecto comercial, se parte del flujo de efectivo, la valuación y el gasto de capital, y se pregunta «cómo se recupera este dinero»; si la IA se ve como la siguiente generación de plataformas, como capacidad de defensa nacional y como una reorganización del orden tecnológico, también hay que preguntar «si no se invierte, ¿se corre el riesgo de quedar desplazado?».

\begin{figure}[H]
\centering
\includegraphics[width=0.96\textwidth]{ai-war-alliances.png}
\caption{Las dos líneas principales de la AI War: el ecosistema de GPU de Nvidia, el ecosistema de TPU de Google y las alianzas de contrapeso que se forman en torno a OpenAI/Google. Diagrama conceptual de elaboración propia, basado en la conversación de 00:15:04--00:17:48.}
\end{figure}

\begin{knowledgebox}{Lectura de la figura: primero los ecosistemas, luego las alianzas}
Arriba a la izquierda y arriba al centro están, respectivamente, las dos líneas principales de ecosistemas de hardware: las GPU de Nvidia y las TPU de Google. OpenAI, Anthropic y otras empresas líderes de modelos se ubican sobre todo en el ecosistema de GPU; Google, en cambio, integra de extremo a extremo modelos, chips, nube y productos. Abajo, la «alianza anti-Google» y la «alianza anti-OpenAI» recuerdan que cuanto más fuerte es un actor individual, más contrapesos provoca por parte de otros gigantes, del capital y de los socios del ecosistema.
\end{knowledgebox}

\subsection{Primer sentido de la AI War: defensa comercial}

Antes se dijo que la AI War no es una figura retórica; este apartado examina primero por qué, en el plano comercial, nadie puede permitirse perder. Guang Mi da dos ejemplos concretos. El primero es que los proveedores de nube quedan reducidos a marca blanca frente a los modelos: antes, desarrolladores y empresas consideraban primero AWS, Azure o Google Cloud; ahora muchas aplicaciones nuevas preguntan primero «qué modelo usar»: GPT, Claude o Gemini. La nube sigue siendo importante, pero la atención de los usuarios podría desplazarse de la nube al modelo.

El segundo es que el Agent podría convertirse en un nuevo punto de entrada del tráfico. La búsqueda tradicional, las Super App, las OTA, las plataformas de transporte, de empleo y de comercio electrónico construyen muchos de sus modelos de negocio sobre la distribución de información, la intermediación y las comisiones por transacción. Si un Agent puede comparar, negociar, hacer pedidos, pagar y dar seguimiento en nombre del usuario, las comisiones y el control del punto de entrada de las plataformas tendrán que valorarse de nuevo. Para los gigantes no se trata de «una herramienta más», sino de una batalla defensiva por fuentes de utilidades que ya tienen.

\begin{warningbox}{Advertencia de interpretación: AI War no significa que las cuentas de corto plazo cuadren}
Decir que la IA es una guerra no implica que cada proyecto de centro de datos, cada plan de entrenamiento de modelos y cada startup merezcan inversión. Solo indica que la conducta de los actores clave no puede explicarse con el ciclo de retorno de una inversión SaaS común, porque el costo de oportunidad de «no invertir» puede ser que se reescriba su posición de plataforma.
\end{warningbox}

\subsection{Segundo sentido de la AI War: estrategia nacional y movilización de capacidad de cómputo}

Guang Mi va más allá y compara la IA con una nueva forma de defensa nacional y con las armas nucleares. Esta frase puede oírse como una exageración, pero apunta a la estructura de capacidades a escala nacional: quien tenga modelos más fuertes, más capacidad de cómputo, una mayor concentración de investigadores y una cadena de suministro de chips más completa podría obtener una ventaja sistémica en la producción de software, el descubrimiento científico, el apoyo militar, la automatización industrial y la distribución de información.

Por eso proyectos cuyas cuentas de flujo de efectivo no cuadran a corto plazo siguen recibiendo apoyo del capital y de las políticas públicas. Las empresas temen quedar desplazadas, los Estados temen rezagarse estratégicamente, y los mercados de capital buscan las pocas direcciones capaces de absorber enormes sumas sin dejar de ofrecer expectativas de crecimiento. Con estos tres factores superpuestos, el ritmo del gasto de capital en IA difícilmente se explica solo con la palabra «burbuja».

\begin{importantbox}{Las tres cuentas para analizar la inversión en IA}
\begin{tabularx}{\textwidth}{p{0.20\textwidth}p{0.35\textwidth}X}
\toprule
Cuenta & Pregunta central & Juicio correspondiente en este episodio \\
\midrule
Cuenta financiera & Cómo cubren la inversión los ingresos, el margen bruto, el flujo de efectivo y la depreciación & El compromiso de 1.4T de OpenAI será cuestionado a corto plazo. \\
Cuenta de plataforma & Si el nuevo punto de entrada sustituirá al anterior y quién controla la relación con el usuario & El Agent y la nube de modelos podrían reescribir la búsqueda, la nube y las Super App. \\
Cuenta estratégica & Si, al no invertir, el país y la empresa pierden capacidades futuras & La IA se compara con la defensa nacional y las armas nucleares; los actores no se atreven a retirarse a la ligera. \\
\bottomrule
\end{tabularx}
\end{importantbox}

\subsection{Del debate sobre la burbuja al análisis del panorama competitivo}

El debate sobre la burbuja suele preguntar «si la valuación es demasiado alta». La pregunta que más vale conservar de este episodio es «quién se verá obligado a seguir invirtiendo, quién puede formar alianzas y quién logra convertir la inversión en un ciclo cerrado de producto y datos». Ese es también el orden de lo que sigue: primero se desglosa la cuenta de ingresos de OpenAI, luego se examinan los bandos de Nvidia/Google, después se discute por qué los tres grandes modelos se alternan el liderazgo y, por último, se llega al Online Learning, una variable que podría cambiar el panorama actual.

\subsection{Resumen de la sección}

El marco de la AI War devuelve la discusión de los ánimos sobre la valuación a las variables estructurales. No niega que haya burbujas parciales, pero recuerda al lector que, cuando la IA amenaza al mismo tiempo a la nube, la búsqueda, las Super App, el trabajo de oficina y la competitividad nacional, los gigantes y los Estados invertirán de forma más agresiva que en un proyecto comercial común.
\section{Las cuentas de ingresos de OpenAI: ingresos visibles e ingresos potenciales no visibles}

La sección anterior explicó por qué no es fácil que los gigantes se levanten de la mesa de juego; esta sección vuelve a enfocar a OpenAI. La mayor controversia en torno a OpenAI es si sus enormes compromisos de capacidad de cómputo corresponden a su flujo de caja futuro. La forma en que Guangmi (广密) desglosa el tema resulta muy didáctica: primero separa los «ingresos visibles» y después los «ingresos no visibles». Los primeros se parecen a los de las plataformas de internet y los servicios en la nube; los segundos dependen de que los Agent puedan sustituir trabajo humano y de que la IA pueda realizar tareas de alto valor que las personas no pueden hacer.

\begin{figure}[H]
\centering
\includegraphics[width=0.92\textwidth]{openai-revenue-stack.png}
\caption{Pila de ingresos de OpenAI: suscripciones, publicidad y comercio electrónico, API, trabajo de los Agent y científicos de IA. Diagrama conceptual de elaboración propia, basado en la conversación de 00:07:38--00:13:10 y 00:13:32--00:15:04.}
\end{figure}

\begin{knowledgebox}{Lectura de la figura: de los «ingresos de plataforma de internet» al «presupuesto de trabajo»}
Las tres primeras capas son ingresos relativamente visibles: suscripciones, publicidad y comercio electrónico, y API. Las dos últimas son más difíciles de estimar: si los Agent pueden sustituir de principio a fin el trabajo de oficina, captan el presupuesto de salarios; si la IA puede escribir CUDA, hacer descubrimientos científicos o buscar nuevos materiales, crea un valor incremental que hoy las personas no han logrado producir plenamente.
\end{knowledgebox}

\subsection{Ingresos visibles: suscripciones, publicidad y comercio electrónico, y API}

La primera categoría de ingresos de Guangmi son las suscripciones. Puede expresarse con una fórmula simplificada:
\[
R_{\text{sub}} = U \times p \times a
\]
donde \(R_{\text{sub}}\) representa el ingreso anual por suscripciones; \(U\), el número de usuarios activos mensuales o de usuarios alcanzables; \(p\), la tasa de conversión a pago; y \(a\), el ingreso anual promedio por usuario de pago. La estimación optimista que se da en el programa es la siguiente: si ChatGPT llega a miles de millones de usuarios activos y alrededor del 10\% paga, los ingresos por suscripciones pueden alcanzar el orden de decenas de miles de millones de dólares. Esta analogía proviene de productos de suscripción consolidados como Office y Netflix.

La segunda categoría son la publicidad y el comercio electrónico. La lógica es que un punto de acceso con mucho tráfico, uso frecuente y alta retención termina por monetizarse. En sus inicios, los videos cortos y la publicidad en teléfonos móviles también se cuestionaron porque la pantalla era pequeña y el contexto no parecía adecuado, pero después se demostró que la publicidad y las transacciones se trasladan adonde está la atención de los usuarios. Si ChatGPT integra búsqueda, recomendaciones, compras e intermediación de servicios, podría quedarse con una parte de la publicidad de búsqueda tradicional y de las comisiones de las plataformas.

La tercera categoría es la API, es decir, la «nube de modelos». Guangmi la compara con el surgimiento de AWS a partir del comercio electrónico de Amazon: cuando una empresa necesita para sí misma enormes sistemas de cómputo, datos e ingeniería, esa misma capacidad puede ofrecerse como servicio a los desarrolladores. Para OpenAI, cuanto más capaz sea el modelo, menor su costo y mayor su ecosistema de desarrolladores, más se parecerán los ingresos de la API a un nuevo negocio en la nube.

\begin{importantbox}{La condición común de los ingresos visibles}
Las suscripciones, la publicidad y el comercio electrónico, y la API dependen de lo mismo: ChatGPT debe mantener mucho tráfico, uso frecuente, una marca fuerte y una experiencia de modelo suficientemente buena. Si el punto de acceso de los usuarios se traslada a Gemini, a los fabricantes de teléfonos o a Agent verticales, el techo de estas tres categorías de ingresos bajará.
\end{importantbox}

\subsection{Ingresos no visibles: Agent que sustituyen trabajo humano}

La expectativa más grande proviene del presupuesto de trabajo. El software de herramientas vende eficiencia; si un Agent realmente completa tareas de principio a fin, lo que vende podría ser «la sustitución de una parte de los salarios». El desarrollo de software es el primer sector en el que aparece esta señal: AI Coding generó en un año un mercado del orden de diez mil millones de dólares, lo que muestra que en el trabajo del conocimiento sí existen tareas que los modelos pueden absorber con rapidez.

Pero aquí hay dos umbrales decisivos. Primero, el Agent debe ser lo bastante confiable para asumir responsabilidades, no solo para generar sugerencias; segundo, las empresas deben estar dispuestas a trasladar presupuesto de personal, subcontratación o software tradicional a sistemas de IA. Solo al superar estos dos umbrales la IA deja de limitarse a «ayudar a los empleados a trabajar más rápido» y se convierte en una nueva oferta de trabajo.

\begin{warningbox}{El presupuesto para herramientas y el presupuesto de salarios no son lo mismo}
Que un empleado esté dispuesto a pagar unas decenas de dólares al mes por una herramienta no significa que la empresa esté dispuesta a entregarle a un Agent el presupuesto de un puesto. Lo segundo exige estabilidad, control de permisos, auditoría, asignación de responsabilidades y resultados predecibles. Muchas empresas de IA se quedarán estancadas en ser «una herramienta muy útil» en lugar de «un sistema capaz de captar el presupuesto de salarios».
\end{warningbox}

\subsection{Ingresos más lejanos: la IA hace lo que las personas no pueden}

Guangmi también menciona otra categoría de ingresos no visibles: la IA podría completar tareas de alto valor que hoy las personas no pueden hacer o hacen muy lentamente, por ejemplo, reescribir CUDA, descubrir nuevos materiales, acelerar la investigación médica, buscar pistas sobre la superconductividad a temperatura ambiente e incluso «que la IA sea científica de IA». Este tipo de ingresos no puede clasificarse simplemente como sustitución de lo existente, porque podría crear nuevos problemas, nuevos mercados y nuevas rutas científicas.

Aquí la condición es más exigente: el modelo no solo debe saber utilizar el conocimiento existente, sino también mejorar mediante experimentación a largo plazo, retroalimentación, razonamiento, memoria y aprendizaje autónomo. Esto conduce de forma natural al Online Learning que se trata más adelante. Dicho de otro modo, cuanto mayor sea la expectativa de ingresos a largo plazo de OpenAI, más altas serán las exigencias de confiabilidad del modelo, aprendizaje continuo y retroalimentación de entornos reales.

\subsection{Resumen de la sección}

Las cuentas de ingresos de OpenAI pueden dividirse en tres niveles: las suscripciones, la publicidad y el comercio electrónico, y la API son ingresos de plataforma visibles; la sustitución del trabajo de oficina por Agent es una migración más grande de valor existente; y los descubrimientos científicos de la IA y las tareas incrementales de alto valor dependen de avances más fundamentales en la capacidad de los modelos. Las dudas a corto plazo provienen de las cuentas financieras; el optimismo a largo plazo, de las cuentas de plataforma y de paradigma.
\section{Dos bandos y tres modelos: por qué el liderazgo alternado se vuelve la norma}

Las dos secciones anteriores trataron la inversión y los ingresos; esta sección entra en el panorama competitivo. La pregunta aquí es: si todos siguen invirtiendo, ¿en torno a qué ejes se organizará la industria? Guangmi (广密) simplifica el ecosistema de hardware de la AI War en dos ejes: el ecosistema de GPU de Nvidia y el ecosistema de TPU de Google; y simplifica el eje de la capa de aplicaciones y modelos en el liderazgo alternado de tres grandes: OpenAI, Anthropic y Google Gemini. Esta simplificación es muy útil, porque evita tomar cada nueva clasificación trimestral como el desenlace final y nos recuerda situar «quién lidera a corto plazo» dentro de una competencia de ecosistemas de más largo plazo.

\subsection{Nvidia GPU vs Google TPU}

Esta sección comienza por el ecosistema de hardware de base, porque la capacidad, el costo y la velocidad de iteración de las empresas de modelos dependen de los chips, la pila de software y la oferta de nube. Solo al entender la diferencia entre GPU y TPU se entiende por qué, aunque en ambos casos se trate de competencia de modelos, Google y Nvidia representan dos formas de organización industrial completamente distintas.

La TPU es la Tensor Processing Unit que Google desarrolló internamente, optimizada para el entrenamiento y la inferencia de aprendizaje automático; la GPU es un ecosistema de cómputo paralelo de propósito general dominado por Nvidia, que en torno a CUDA, los desarrolladores, los servicios de nube, los servidores, las redes y la cadena de suministro ha generado poderosos efectos de red. La ventaja de Google es la integración de extremo a extremo: modelos, chips, nube, búsqueda, Android, YouTube y Workspace están dentro de una sola empresa; la ventaja de Nvidia es el ecosistema abierto: OpenAI, Anthropic, xAI, Meta, los proveedores de nube y las startups pueden formar una red de innovación externa en torno a la GPU.

Guangmi compara a Google con el Apple de la era de la IA y a Nvidia con el Android de la era de la IA. Esto no significa que ambos sean completamente equivalentes, sino que señala dos formas de organización: la integración vertical y el ecosistema abierto. La primera puede cerrar el ciclo en producto y costo; la segunda es más flexible en cuanto a desarrolladores, innovación entre múltiples empresas y expansión de la oferta.

\begin{knowledgebox}{Términos clave: GPU, TPU y CUDA}
\begin{tabularx}{\textwidth}{p{0.18\textwidth}p{0.32\textwidth}X}
\toprule
Término & Problema que resuelve & Relación con este episodio \\
\midrule
GPU & Acelera las operaciones matriciales y el aprendizaje profundo mediante cómputo paralelo a gran escala & Núcleo del ecosistema de Nvidia; sostiene a empresas de modelos externas como OpenAI/Anthropic. \\
TPU & Chip acelerador que Google diseñó a la medida para el aprendizaje automático & Sostiene el ciclo vertical cerrado de Google, desde el chip hasta Gemini. \\
CUDA & Ecosistema de software de cómputo paralelo de Nvidia & Hace que la GPU no sea solo hardware, sino una cadena de herramientas para desarrolladores e ingeniería. \\
Nicho ecológico & Posición estable de una empresa en la cadena industrial & Google hace full-stack, Nvidia construye el ecosistema de base y las empresas de modelos aportan capacidades y productos. \\
\bottomrule
\end{tabularx}
\end{knowledgebox}

\subsection{El liderazgo alternado de GPT, Claude y Gemini}

El liderazgo alternado de los tres grandes modelos es uno de los juicios sobre el panorama competitivo más importantes de este episodio. Guangmi considera que, bajo el paradigma actual de preentrenamiento más RL/Post-training, los métodos, el talento y la información de los equipos de punta circulan muy rápido, por lo que es difícil que una sola empresa mantenga una gran ventaja durante mucho tiempo. Así, la competencia pasa de «quién aplasta a los demás de una sola vez» a «quién optimiza de extremo a extremo con mayor profundidad una Beta estratégica determinada».

\begin{figure}[H]
\centering
\includegraphics[width=0.96\textwidth]{top-model-rotation.png}
\caption{Liderazgo alternado de los tres grandes modelos: GPT, Claude y Gemini tienen prioridades estratégicas distintas. Diagrama conceptual de elaboración propia, basado en la conversación de 00:17:48--00:25:40.}
\end{figure}

\begin{knowledgebox}{Lectura de la figura: el liderazgo alternado no significa ausencia de diferenciación}
GPT es más fuerte en productos para el consumidor, marca y posicionamiento como asistente personal; Claude es más fuerte en Coding, empresas y la experiencia de los trabajadores del conocimiento; Gemini es más fuerte en multimodalidad, el ecosistema de Google y la distribución a través de Android y la búsqueda. El liderazgo alternado muestra que es difícil sostener a largo plazo una ventaja en capacidades generales, pero la diferenciación muestra que cada empresa optimizará en torno a sus propios datos y objetivos de producto.
\end{knowledgebox}

\begin{warningbox}{La lectura equivocada del SOTA de los benchmarks}
SOTA significa originalmente state of the art, es decir, el nivel más avanzado del momento. Pero este episodio subraya que el SOTA de una clasificación no equivale al SOTA comercial. Los indicadores realmente importantes pueden ser los ingresos, los clientes, los usuarios activos diarios, la retención, la producción de tokens, la tasa de éxito en tareas y cuánto trabajo real están dispuestos a delegar los usuarios.
\end{warningbox}

\subsection{Por qué las empresas de modelos fundacionales se parecen al comercio electrónico generalista}

Guangmi propone una analogía «perezosa pero útil»: las empresas de modelos fundacionales se parecen al comercio electrónico generalista de hace diez años; Scale SKU es análogo a Scale Data. Cuando el comercio electrónico generalista se expandió de los libros a todas las categorías, pudo reutilizar almacenes, logística, pagos, tráfico y relación con los usuarios; cuando las empresas de modelos fundacionales se expanden del chat a Coding, Office, investigación de inversiones, medicina, diseño y multimodalidad, también pueden reutilizar la arquitectura del modelo, los pipelines de datos, los sistemas de evaluación, los procesos de RL y los puntos de acceso de los usuarios.

\begin{figure}[H]
\centering
\includegraphics[width=0.94\textwidth]{foundation-model-as-ecommerce.png}
\caption{Modelo fundacional = comercio electrónico generalista: Scale SKU es análogo a Scale Data. Diagrama conceptual de elaboración propia, basado en la conversación de 00:25:40--00:27:40.}
\end{figure}

\begin{knowledgebox}{Lectura de la figura: el centro de la analogía es «reutilizar la infraestructura»}
El comercio electrónico generalista pudo pasar de los libros a todas las categorías no solo porque quería vender más cosas, sino porque el almacenamiento y la distribución, los pagos, las recomendaciones y la confianza de los usuarios podían reutilizarse. Las empresas de modelos fundacionales pueden pasar del chat a Coding, Office y Agents sectoriales no simplemente agregando funciones, sino porque los pipelines de recolección de datos, entrenamiento, evaluación y retroalimentación de producto pueden reutilizarse.
\end{knowledgebox}

Esta analogía deja dos advertencias para los emprendedores. Primera: los escenarios verticales no necesariamente carecen de valor, pero hay que saber si uno está construyendo un «comercio electrónico vertical» o un «Xiaohongshu (小红书)». El comercio electrónico vertical puede ser rentable, pero el techo de la plataforma es limitado; Xiaohongshu construyó, al margen del comercio electrónico generalista, activos de contenido y un posicionamiento comunitario, y por eso tiene un valor distinto. Segunda: las empresas de modelos son más eficientes al bajar hacia los datos, así que las startups deben encontrar datos diferenciados, flujos de trabajo o relaciones con los usuarios que a las empresas de modelos fundacionales no les resulte fácil obtener.

\subsection{Resumen de la sección}

Esta sección descompuso el panorama competitivo en ecosistemas de hardware, los tres grandes modelos y la analogía de plataforma. El liderazgo alternado muestra que, bajo el paradigma actual, la brecha tecnológica difícilmente se vuelve permanente; la Beta estratégica y el ciclo cerrado de datos determinan dónde cada empresa construye su diferenciación. Que los modelos fundacionales se parezcan al comercio electrónico generalista significa que su capacidad de «expandir categorías horizontalmente» es muy grande, y las startups deben evitar las posiciones que pueden absorberse fácilmente mediante Scale Data.
\section{Google, OpenAI y Search: cómo se redistribuye el control de la puerta de entrada}

La sección anterior mostró que Google y OpenAI están en la mesa principal; esta sección analiza específicamente la competencia entre ambos por la puerta de entrada. El juicio más tenso del programa es este: el regreso triunfal de Google no significa que la crisis haya pasado, y ChatGPT tampoco es solo una herramienta de chat marginal. Al final, ambos podrían enfrentarse directamente en la búsqueda, los asistentes personales, los dispositivos móviles y los presupuestos publicitarios.

\subsection{Tras el ascenso de Gemini, ¿por qué OpenAI sigue siendo fuerte?}

Guangmi (广密) tiene una visión más bien optimista de OpenAI, por tres razones. Primera: OpenAI podría ser una de las pocas empresas de los últimos veinte años con una posibilidad real de desafiar a Google. Ni siquiera Bing, de Microsoft, cambió de fondo la puerta de entrada de la búsqueda; ChatGPT, en cambio, trasladó la mentalidad del usuario de «buscar respuestas» a «que la IA me ayude a entender, generar y ejecutar». Segunda: OpenAI sigue siendo uno de los mejores activos de IA: su densidad de talento, su sentido de producto, su marca y su escala de usuarios son muy sólidos. Tercera: OpenAI conserva una probabilidad alta de participar en la exploración del siguiente paradigma, lo cual es decisivo para su valuación a largo plazo.

Las ventajas de Google son igual de reales: activos como Gemini, la multimodalidad, Android, Search, YouTube, Workspace, Cloud, las TPU y Waymo forman una base muy amplia. El problema es que la fortaleza de Google no elimina automáticamente a OpenAI; este podría ser un mercado donde coexistan a largo plazo dos o incluso varios gigantes.

\subsection{Cómo reescribirá el Agent a Search}

Si, como se dijo antes, la crisis de Google no ha pasado realmente, la cuestión recae en Search: la búsqueda no es una función aislada, sino la puerta de entrada general a la publicidad, la distribución de páginas web, el tráfico hacia los comercios y las decisiones de los usuarios. Esta sección pone lado a lado el flujo de la búsqueda tradicional y el flujo de un Agent; solo así se ve por qué la redistribución del tráfico toca el núcleo de las utilidades de Google.

El flujo tradicional de la búsqueda es el siguiente: el usuario escribe palabras clave, el motor de búsqueda devuelve enlaces, resúmenes y anuncios, y el usuario hace clic, compara, compra o sigue buscando. El flujo de un Agent podría ser, en cambio: el usuario plantea un objetivo y el Agent busca, compara, invoca herramientas, negocia el precio, hace el pedido, paga, da seguimiento y, al final, devuelve el resultado. La diferencia entre ambos no es solo de interfaz, sino de posición en la cadena de valor.

\begin{figure}[H]
\centering
\includegraphics[width=0.96\textwidth]{search-agent-disruption.png}
\caption{Redistribución del tráfico entre Search y el Agent: de los enlaces de búsqueda a la tarea completada. Diagrama conceptual propio, elaborado a partir de la conversación de 00:31:20--00:35:08 y 00:50:57--00:52:53.}
\end{figure}

\begin{knowledgebox}{Lectura de la figura: lo que el Agent disputa es el poder de distribución}
En la figura, el objetivo del usuario entra primero al Agent, y no directamente al cuadro de búsqueda ni a una Super App. Después, el Agent decide si invoca la búsqueda, la venta de boletos, el transporte, los pagos o los sistemas empresariales. Quien controle esta puerta de entrada tendrá la oportunidad de redistribuir la publicidad, las transacciones, las comisiones y los datos de los usuarios.
\end{knowledgebox}

\begin{warningbox}{El Agent no eliminará la búsqueda de un momento a otro}
Para que el Agent reescriba de verdad la publicidad en buscadores, debe resolver la confiabilidad de las citas, la información en tiempo real, los pagos, los permisos, la integración de los comercios, la prevención de fraudes, la auditoría y la asignación de responsabilidades. A corto plazo es más probable una forma híbrida: la búsqueda tradicional incorpora IA, ChatGPT incorpora la búsqueda, y cada uno va erosionando el terreno del otro.
\end{warningbox}

\subsection{La mentalidad móvil de ChatGPT y la mentalidad de ecosistema de Gemini}

A continuación hay que pasar del mecanismo de la puerta de entrada a la mentalidad del usuario, porque modelos que responden igual de bien pueden terminar abriéndose en momentos, dispositivos y escenarios completamente distintos. Esta sección trata de la frecuencia de uso y del lugar que ocupa cada uno en la vida diaria, no solo de la capacidad del modelo.

En el programa se menciona una observación con un fuerte enfoque de producto: ChatGPT se parece más a un asistente para la vida diaria, que muchos usuarios usan los fines de semana, en el celular y en lo cotidiano; Gemini, en algunos escenarios, se parece más a una herramienta de productividad para el horario laboral, y cuenta con una ventaja natural de distribución gracias a Android y a los mercados emergentes. Esta diferencia implica que, además de la capacidad del modelo, la frecuencia de uso, la cercanía del escenario y la mentalidad del usuario importan igual.

Si ChatGPT logra convertirse en un asistente personal con 1,000--2,000 millones de usuarios activos diarios, su base comercial será muy estable; si Gemini logra integrar de verdad Android, Search, Workspace y la multimodalidad, Google también podría mantener su fortaleza en la era de la IA. Los dos caminos no se excluyen entre sí: compiten por cuál puerta de entrada abre primero el usuario.

\subsection{Resumen de la sección}

La crisis de Google pasó de «¿se le escapará la IA?» a «¿podrá conservar el poder de distribución en la puerta de entrada del Agent?». El reto de OpenAI pasó de «¿podrá construir buenos modelos?» a «¿podrá convertir ChatGPT en un asistente personal y una puerta de entrada de búsqueda y ejecución de uso frecuente, confiable y monetizable?». La conclusión de esta sección es que la guerra por la puerta de entrada apenas comienza, y los rankings de modelos son solo una pequeña parte de ella.
\section{El tercer paradigma, Online Learning: por qué podría ser más fundamental que los robots, los modelos del mundo y la multimodalidad}

Lo anterior trató el panorama competitivo actual; esta sección entra en los paradigmas técnicos. Guangmi (广密) considera que el scaling del preentrenamiento está por terminar y que RL/Post-training todavía tiene margen, pero que lo que realmente podría cambiar el panorama en la siguiente etapa es el Online Learning o Continuous Learning. Lo central aquí no se reduce a que «el modelo tenga un contexto más largo», sino a si el modelo puede aprender de forma continua a partir de la interacción, la retroalimentación y las tareas reales.

\begin{figure}[H]
\centering
\includegraphics[width=0.96\textwidth]{scaling-paradigm-shift.png}
\caption{Tres paradigmas de aprendizaje: preentrenamiento, RL/Post-training, Memory y Online Learning. Diagrama conceptual propio, elaborado a partir de la conversación de 00:36:01--00:41:01.}
\end{figure}

\begin{knowledgebox}{Lectura de la figura: de la compresión estática al aprendizaje por interacción}
El preentrenamiento equivale a comprimir de una sola vez datos de internet y datos especializados a gran escala; RL/Post-training moldea el modelo con retroalimentación humana, procesos de expertos y resultados de tareas; Memory y el contexto largo permiten que el modelo recuerde más historial del usuario y de las tareas; Online Learning va más allá y exige que el modelo se actualice realmente durante la inferencia y la interacción, de modo que entienda mejor al usuario cuanto más se usa y haga mejor las tareas cuanto más las realiza.
\end{knowledgebox}

\subsection{Preentrenamiento: un compresor gigantesco}

El preentrenamiento (Pre-training) consiste en entrenar un modelo con cantidades masivas de texto, código y datos multimodales para que aprenda a predecir el siguiente token o la siguiente estructura. Guangmi enfatiza que los modelos actuales siguen siendo compresores gigantescos: las tareas que no están en la distribución de los datos suelen no funcionar, y el modelo solo es más confiable cuando ha comprimido datos similares.

Esta afirmación es fundamental. Explica por qué los modelos ya superan a la mayoría de las personas en conocimiento y, aun así, siguen fallando en flujos de trabajo reales. El trabajo real no consiste solo en responder preguntas, sino en operar dentro de un CRM, sistemas financieros, software de diseño, procesos hospitalarios, ventas en campo, permisos empresariales y cadenas de herramientas concretas. Si el modelo no ha visto suficientes procesos reales, solo puede recurrir a conjeturas por generalización.

\begin{importantbox}{La perspectiva del compresor}
El modelo no «entiende» una tarea de la nada, sino que comprime los patrones de los datos de entrenamiento en parámetros y representaciones. Si una tarea nunca ha entrado en la distribución de los datos, al modelo le resulta difícil completarla de forma estable; para que un Agent sea confiable, hay que convertir los procesos de trabajo reales en datos de los que se pueda aprender.
\end{importantbox}

\subsection{RL/Post-training: datos de expertos y retroalimentación de tareas}

Antes se explicó el preentrenamiento como una compresión a gran escala; esta subsección examina «cómo se moldea el modelo después de la compresión». Si el preentrenamiento permite que el modelo aprenda la gran cantidad de patrones que ya existen en el mundo, RL/Post-training se ocupa más de cómo el modelo se vuelve útil en tareas concretas, ante las preferencias del usuario y en los procesos de expertos.

RL es Reinforcement Learning, aprendizaje por refuerzo. No se reduce a las palabras «modelo de recompensa», sino que es una forma de entrenamiento en la que el sistema ajusta su comportamiento según la retroalimentación de los resultados. En el contexto de los modelos grandes, Post-training suele incluir ajuste por instrucciones, retroalimentación de preferencias humanas, trayectorias de expertos, evaluación automática, retroalimentación sobre el uso de herramientas y calificación de resultados de tareas. En el programa se compara este tipo de datos con las «nuevas energías»: son útiles, pero su volumen total es reducido y su recolección es costosa.

Para que un Coding Agent mejore, se necesitan código de alta calidad, issues reales, pruebas, retroalimentación de PR y procesos de ingeniería de software de largo plazo; para que un Agent médico mejore, se necesitan los procesos de diagnóstico de los médicos, expedientes clínicos, estudios, medicación, seguimiento y límites de responsabilidad; para que un Agent de análisis de inversiones mejore, se necesitan los flujos de trabajo de los analistas, las fuentes de datos, los supuestos de los modelos y la revisión de errores. Ninguna de estas categorías es texto web común: todas son datos de procesos de expertos.

\begin{knowledgebox}{Términos clave: vocabulario de RL/Post-training}
\begin{tabularx}{\textwidth}{p{0.22\textwidth}p{0.30\textwidth}X}
\toprule
Término & Mecanismo central & Papel en este episodio \\
\midrule
RL & Optimiza la política con recompensas o retroalimentación de resultados & Hace que el modelo no solo imite texto, sino que se optimice hacia el éxito en la tarea. \\
Post-training & Alineación, instrucciones, preferencias y entrenamiento por tareas posteriores al preentrenamiento & Determina si el modelo se ajusta a los objetivos del producto. \\
Trayectorias de expertos & Registro del proceso con el que expertos de primer nivel completan una tarea & Permite que el modelo aprenda flujos de trabajo reales, no solo respuestas. \\
Credit assignment & Atribuir el éxito o el fracaso a los pasos intermedios & En la sección de Robotics se usa para explicar si cada paso favorece el éxito de la tarea. \\
Value function & Función que estima el valor futuro del estado o la acción actuales & Ayuda al sistema a saber qué pasos están más cerca del éxito. \\
\bottomrule
\end{tabularx}
\end{knowledgebox}

\subsection{Online Learning: aprender mientras se interactúa}

En este episodio, Online Learning se refiere a que el modelo aprenda de manera autónoma y continua, y que absorba la experiencia del usuario y de las tareas a partir de la interacción. A diferencia del preentrenamiento estático, exige que el modelo, a medida que se usa, entienda mejor al usuario y la tarea y actúe con más iniciativa. El escenario de producto que plantea Guangmi es este: hoy ChatGPT tiene que esperar el prompt del usuario; en el futuro, la IA podría reflexionar, prepararse y establecer asociaciones en segundo plano, de modo que, cuando el usuario vuelva a preguntar, ya tenga una comprensión más profunda.

Si esta capacidad se materializa, cambiará la forma de los productos. La IA dejará de ser una herramienta que razona desde cero cada vez y pasará a ser un compañero de largo plazo que crece junto con el usuario; dejará de limitarse a responder la pregunta actual y recordará los objetivos, preferencias, proyectos, contexto y errores pasados del usuario; dejará de esperar órdenes de forma pasiva y se preparará y emitirá recordatorios por iniciativa propia.

\begin{importantbox}{Definición de Online Learning}
En estos apuntes, Online Learning designa la capacidad de un modelo o sistema de IA de aprender de la interacción continua, de los resultados de las tareas y de la retroalimentación del usuario, y de aplicar lo aprendido de forma estable a su comportamiento futuro. No se trata de guardar el historial de conversaciones, sino de convertir la experiencia en actualizaciones de capacidades, preferencias o estrategias.
\end{importantbox}

\begin{warningbox}{Memory no equivale a Online Learning}
Memory puede limitarse a almacenar información del usuario; el contexto largo puede limitarse a meter más historial en el prompt. Online Learning exige que el sistema aprenda de ese historial mejoras de comportamiento generalizables. Saber «recordar hechos» no equivale a «volverse más inteligente cuanto más se usa».
\end{warningbox}

\subsection{La metáfora energética: petróleo, nuevas energías y fusión nuclear}

Esta subsección usa una metáfora energética para condensar las tres etapas de aprendizaje anteriores en un marco comparativo fácil de recordar. El valor de la metáfora no está en predecir con precisión la ruta tecnológica, sino en ayudar al lector a ver a la vez las diferencias de restricciones entre tres tipos de recursos: los datos existentes, los datos de expertos y el aprendizaje continuo.

Guangmi usa tres metáforas energéticas para explicar las diferencias entre paradigmas: los datos de preentrenamiento son como el petróleo, abundantes pero finitos; los datos de expertos para RL son como las nuevas energías, útiles pero escasos en total; Online Learning es como la fusión nuclear, que todavía no logra un avance real, pero que, si lo logra, tendrá una densidad energética y un techo completamente distintos.

\begin{figure}[H]
\centering
\includegraphics[width=0.96\textwidth]{energy-metaphor.png}
\caption{Metáfora energética: el petróleo, las nuevas energías y la fusión nuclear corresponden a tres tipos de datos y de formas de aprendizaje. Diagrama conceptual propio, elaborado a partir de la conversación de 00:41:01--00:43:05.}
\end{figure}

\begin{knowledgebox}{Lectura de la figura: los límites de la metáfora}
El petróleo subraya que los datos existentes son finitos; las nuevas energías, que los datos de expertos son útiles pero difíciles de recolectar; la fusión nuclear, el techo potencial de Online Learning. La metáfora no implica que la tecnología vaya a repetir necesariamente la historia de la energía; solo nos ayuda a distinguir las restricciones de tres tipos de datos y mecanismos de aprendizaje.
\end{knowledgebox}

\subsection{Por qué los robots, los modelos del mundo y la multimodalidad podrían ser «falsos problemas»}

La afirmación más polémica del programa es esta: muchos de los temas de robots, modelos del mundo y multimodalidad podrían ser falsos problemas, y Online Learning sería el verdadero problema. Aquí «falso problema» no significa que estas líneas no sean importantes, sino que podrían depender de una generalización y un aprendizaje continuo más fundamentales. Sin generalización, los robots acabarían como la conducción autónoma: cada escenario de cola larga requiere recolectar datos, etiquetarlos y ajustar reglas, y se cae por mucho tiempo en la situación de «tanta inteligencia como trabajo humano haya».

Con la multimodalidad y los modelos del mundo ocurre algo parecido. Que un modelo pueda ver video, generar imágenes y predecir procesos físicos no significa que pueda aprender de forma continua en entornos reales, transferir experiencia y resolver tareas nuevas. Lo verdaderamente difícil es convertir la experiencia en capacidades reutilizables. Si Online Learning logra un avance, los robots, los modelos del mundo y la multimodalidad podrían obtener una generalización mucho mayor; si no, estas líneas seguirán frenadas por el volumen de datos, el costo del etiquetado y la complejidad de los entornos reales.

\subsection{Resumen de la sección}

Lo central del tercer paradigma no es un término nuevo, sino el paso de la compresión estática al aprendizaje continuo por interacción. El preentrenamiento resuelve la compresión de conocimiento a gran escala; RL/Post-training resuelve la retroalimentación de tareas y los procesos de expertos; Online Learning intenta resolver el problema de «volverse más inteligente cuanto más se usa» y de «generalizar a partir de entornos reales». Su importancia radica en que podría alterar el equilibrio actual, en el que los tres grandes actores se alternan el liderazgo.
\section{El modelo es el producto, los datos son el modelo: Agentic Web y la línea principal de inversión}

La sección anterior señaló que el Online Learning requiere interacción y retroalimentación reales; esta sección lleva esa idea al producto, a los datos y a la estrategia de inversión. Guangmi (广密) insiste en que «el modelo es el producto, los datos son el modelo». La frase parece enrevesada, pero es muy concreta: la mayor mejora en la experiencia de los últimos años vino de la mejora de los modelos, y la raíz de la mejora de los modelos está, a su vez, en la mejora de la distribución de los datos. Quien consiga datos diferenciados podrá diferenciarse en el modelo y en el producto.

\begin{figure}[H]
\centering
\includegraphics[width=0.96\textwidth]{model-data-product-loop.png}
\caption{El modelo es el producto, los datos son el modelo: las tareas de los usuarios, los datos de interacción, la destilación de expertos, la mejora del modelo y la experiencia de producto forman un ciclo cerrado. Diagrama conceptual propio, elaborado a partir de la conversación de 00:43:05--00:56:45.}
\end{figure}

\begin{knowledgebox}{Lectura de la figura: la experiencia de producto vuelve a la distribución de los datos}
Las tareas de los usuarios exponen flujos de trabajo reales; los datos de interacción registran preferencias y puntos de falla; la destilación de expertos aporta procesos de alta calidad; la mejora del modelo amplía el límite de sus capacidades, y una mejor experiencia de producto atrae, a su vez, más tareas y más retroalimentación. Este ciclo cerrado es el sentido práctico de «el modelo es el producto, los datos son el modelo».
\end{knowledgebox}

\subsection{Un eje horizontal y uno vertical: destilación del conocimiento experto y Online Learning}

Guangmi resume el camino futuro como «un eje horizontal y uno vertical». El horizontal consiste en destilar el conocimiento de expertos humanos e incorporar más industrias, tareas y cadenas de herramientas a la distribución de datos del modelo; el vertical abarca el Online Learning, la Memory, los Agent proactivos y nuevas formas de interacción. El eje horizontal permite que la IA cubra más trabajo de oficina; el vertical podría permitir que la IA crezca de manera continua a partir de los usuarios y del entorno.

\begin{importantbox}{Un eje horizontal y uno vertical}
El eje horizontal responde a «qué industrias y tareas todavía no domina», y amplía el límite de capacidades mediante datos de expertos, flujos de trabajo reales y escenarios de cada industria; el eje vertical responde a «si puede volverse más inteligente cuanto más se usa», y se apoya en la Memory, el Online Learning, los Agent proactivos y la interacción a largo plazo.
\end{importantbox}

Esto tiene implicaciones directas para las empresas emergentes. Una UI delgada es fácil de copiar para las empresas de modelos base o las grandes plataformas; pero si se obtienen datos únicos, procesos de expertos, retroalimentación de usuarios y flujos de trabajo de una industria, es posible formar «activos de contenido» o «activos de proceso» sobre el modelo base. Por eso este episodio vuelve una y otra vez a dominios concretos como Coding, Office, Finance, salud, materiales y robótica.

\subsection{Agentic Web: ¿la verdadera Web3?}

Otro concepto importante del programa es la Agentic Web. Aquí Web3 no se refiere a la Web3 en el sentido de blockchain, sino a que, cuando el Agent se convierte en el nuevo punto de entrada, se reorganiza la relación entre las páginas web, las plataformas, los pagos, la búsqueda y las tareas de los usuarios. Antes, los puntos de entrada centrales de internet eran el cuadro de búsqueda, la App, el Feed y la Super App; el punto de entrada central de la Agentic Web podría ser «quiero lograr algo».

Si un Agent puede operar de forma continua, deja de ser solo una herramienta de consumo de contenido y se convierte en el intermediario que distribuye la atención, los pedidos, el tráfico y las transacciones. Para Uber, Ctrip (携程), la publicidad en buscadores, los fabricantes de teléfonos y las superaplicaciones, esta es una variable que obliga a renegociar quién controla el punto de entrada.

\begin{warningbox}{Las carencias de infraestructura de la Agentic Web}
La Agentic Web no se materializa solo con la capacidad del modelo. También necesita identidad, autorización, pagos, control de riesgos, verificación de humanos frente a máquinas, protocolos de sitios web, registros de auditoría y asignación de responsabilidades. Si estos componentes no avanzan al mismo ritmo, el Agent quedará atorado en el punto de «se puede demostrar, pero es difícil ejecutarlo a escala».
\end{warningbox}

\subsection{Estrategia de inversión: apostar donde el crecimiento tecnológico es más pronunciado}

Una vez enlazados el producto, los datos y la Agentic Web en un ciclo cerrado, esta sección pasa al marco de inversión. Aquí se habla de inversión no para dar recomendaciones de compra o venta, sino para usar la clasificación de un inversionista como vía inversa para entender la estructura de la industria: hacia qué capa fluye el dinero suele indicar qué capa es la más escasa y la que más amplifica el cambio tecnológico.

Guangmi resume la estrategia de inversión en IA así: invertir en los modelos más avanzados, invertir en el cómputo y la infraestructura de silicio que esos modelos necesitan, e invertir en los beneficios del derrame tecnológico de esos modelos. Este marco corresponde a tres clases de activos: empresas de modelos y de puntos de entrada como OpenAI/Google/Anthropic/ByteDance (字节); la oferta de base como Nvidia/TSMC/infraestructura de cómputo, y empresas que aprovechan el derrame tecnológico, como Cursor, Perplexity, Coding Agent, etiquetado de datos, Serving y herramientas multimodales. Aquí Perplexity es el nombre de una empresa de búsqueda y preguntas y respuestas con IA, no la perplexity/PPL de la evaluación de modelos de lenguaje; esta última es la métrica de «perplejidad», el exponencial de la entropía cruzada, que representa de forma aproximada la incertidumbre promedio del modelo ante el siguiente token.

El programa también menciona un cambio de cartera: de una apuesta más concentrada en OpenAI y ByteDance a una canasta formada por OpenAI, ByteDance, Google, Anthropic, Nvidia, TSMC y otras. La lógica detrás es la alternancia en el liderazgo y la incertidumbre sobre el paradigma: si es difícil prever al ganador final, conviene armar una cartera en torno a la línea principal en lugar de apostar por una sola empresa.

\begin{knowledgebox}{No es una recomendación de inversión, sino un mapa de la industria}
Estos apuntes conservan el marco de inversión del programa para entender la estructura de la industria: cómo se impulsan mutuamente los modelos, el cómputo, los puntos de entrada de las plataformas, el derrame hacia las aplicaciones y los servicios de datos. No es una recomendación sobre ningún valor ni sobre la valuación de ninguna empresa.
\end{knowledgebox}

\subsection{Los temas de 2026: Online Learning, trabajadores del conocimiento, multimodalidad y Agentic Web}

Esta sección reúne los juicios técnicos e industriales anteriores en una lista de observación para 2026. En lugar de preguntar «quién ganará al final», la pregunta más operativa es: qué temas mostrarán primero señales verificables y qué escenarios generarán primero ingresos y cambios en el comportamiento de los usuarios.

De cara a 2026, el programa enumera varios temas esperados. Primero, el Online Learning o señales del siguiente paradigma; segundo, Agent más proactivos y una Memory más sólida; tercero, capacidades L3/L4 parciales orientadas a los trabajadores del conocimiento, como Coding, Office/PPT/Excel, análisis financiero de inversiones y la obtención de información de cola larga; cuarto, un gran año para la multimodalidad; quinto, el impacto de la Agentic Web sobre el poder de distribuir el tráfico.

Aquí L3/L4 es una analogía tomada de los niveles de conducción autónoma: L3 puede entenderse como que el sistema asume la tarea principal en condiciones específicas, mientras la persona todavía debe supervisar o retomar el control; L4 significa completar la tarea con un alto grado de automatización dentro de escenarios acotados. Guangmi considera que el «L3/L4 parcial» es muy realista, mientras que el «L4 total» sigue siendo muy difícil. Dicho de otro modo, primero se obtiene valor económico en flujos de trabajo bien definidos, en lugar de esperar una AGI capaz de todo.

\subsection{Resumen de la sección}

«El modelo es el producto, los datos son el modelo» enlaza la tecnología, el producto y la inversión en un ciclo cerrado. La experiencia del modelo proviene de los datos, los datos provienen de tareas reales, las tareas reales provienen de los puntos de entrada del producto, y estos, a su vez, traen más retroalimentación. Lo decisivo en 2026 no es quién se declara más general, sino quién logra que los datos y el producto se refuercen de forma acumulativa en escenarios concretos.
\section{Neo Labs y Robotics: nuevos laboratorios, robots y datos del mundo real}

La sección anterior trató el ciclo cerrado de datos; esta examina los nuevos laboratorios de Silicon Valley y Robotics. Guangmi (广密) menciona Neo Labs como SSI, Thinking Machines, Periodic Labs, Aethera, General Intuition y Reflection, y también habla de empresas de robótica como Pi, Generalist, Sundee, 1X, Figure, Google y Tesla. La palabra clave común de esta sección sigue siendo los datos: el siguiente paradigma necesita datos, y los robots necesitan todavía más datos del mundo real.

\begin{figure}[H]
\centering
\includegraphics[width=0.88\textwidth]{neo-labs-map.png}
\caption{Distribución de los Neo Labs: aprendizaje autónomo, Post-training, modelos de materiales, Multi-Agent, datos para modelos del mundo y modelos de código abierto. Diagrama conceptual propio, elaborado a partir de la conversación entre 00:56:45 y 00:59:43.}
\end{figure}

\begin{knowledgebox}{Lectura de la figura: los nuevos laboratorios no repiten la misma apuesta}
SSI parece apostar por el aprendizaje autónomo; a Thinking Machines se le considera un equipo de Post-training y del siguiente paradigma; Periodic Labs se centra en modelos de materiales; Aethera explora entornos Multi-agent; General Intuition construye modelos del mundo con datos de momentos destacados de videojuegos; Reflection se describe como el actor que cubre el hueco del ecosistema estadounidense de modelos de código abierto. Lo que tienen en común es que, tras la salida de talento de OpenAI, el siguiente paradigma recibe apuestas por varias rutas a la vez.
\end{knowledgebox}

\subsection{Por qué OpenAI sigue siendo fuerte después de sus escisiones}

Esta subsección vuelve al tema de la organización del talento, porque la aparición de los Neo Labs no es una ola casual de emprendimiento, sino el resultado de que el talento, las rutas técnicas y la experiencia de las empresas de modelos de frontera se desbordaron hacia fuera. Solo entendiendo de dónde vienen estos equipos se puede juzgar por qué quizá no sean simplemente «otra empresa de modelos más».

El programa ofrece una observación muy interesante: aunque de OpenAI se escindieron Anthropic, SSI de Ilya y los equipos vinculados a Mira y John Schulman, OpenAI sigue siendo muy fuerte. Se considera que Anthropic es una extensión del primer Scaling team de OpenAI; SSI está más cerca de la línea de Pre-training y aprendizaje autónomo, y Thinking Machines, de la experiencia de ChatGPT y Post-training. Esto muestra que la densidad de talento y la exploración de paradigmas de OpenAI llegaron a estar tan concentradas que, al dispersarse, dieron lugar a varios centros nuevos.

Sin embargo, esto también significa que la siguiente etapa no tiene por qué quedar en manos exclusivas de las grandes empresas actuales. Los nuevos laboratorios pueden saltarse la carrera por alcanzar a los modelos existentes y apostar directamente por la siguiente generación de métodos de aprendizaje, la ciencia de materiales, los entornos Multi-Agent o los datos para modelos del mundo. Por eso este episodio subraya tanto el liderazgo de OpenAI como la necesidad de no fijarse solo en las clasificaciones actuales.

\subsection{Robotics: la etapa entre GPT-0.5 y GPT-1}

El juicio central de la parte de robótica es este: el avance es rápido, pero el momento GPT-4 todavía está lejos. Guangmi considera que hoy la robótica podría estar entre GPT-0.5 y GPT-1, y que en los próximos dos o tres años podría llegar a GPT-2 o GPT-3. La razón es que la robótica no sigue el camino de los modelos de lenguaje, que primero se unificaron y después se diversificaron. Los modelos de lenguaje cuentan con token, corpus de texto, Transformer y una ruta de preentrenamiento relativamente unificados; la robótica, en cambio, está fragmentada desde el inicio: hardware distinto, sensores distintos, escenarios distintos, métodos de recolección distintos y objetivos de evaluación distintos.

Hoy la diferenciación entre empresas de robótica se manifiesta en la recolección de datos. Google y Pi tienen datos de robots reales obtenidos por teleoperación; Sundee recolecta datos domésticos con guantes y mochilas; Generalist afirma contar con datos de robots reales a gran escala; los modelos del mundo y los datos de YouTube y de videojuegos también empiezan a entrar en el entrenamiento. Cada ruta responde a la misma pregunta: cómo conseguir que los robots obtengan datos de manipulación en cantidad suficiente, suficientemente reales y de calidad suficientemente alta.

\begin{knowledgebox}{Términos clave: datos robóticos y generalización}
\begin{tabularx}{\textwidth}{p{0.22\textwidth}p{0.30\textwidth}X}
\toprule
Término & Explicación en una frase & Relación con este episodio \\
\midrule
Teleoperación & Una persona controla un robot a distancia para completar una tarea y se registra el proceso & Aporta trayectorias de acción reales y retroalimentación del entorno. \\
Datos de robots reales & Datos de sensores y acciones de robots reales en entornos reales & Más cercanos a la implementación que la simulación, pero costosos de recolectar. \\
Transferencia entre cuerpos robóticos & Compartir o transferir capacidades entre distintos tipos de hardware robótico & Resuelve la fragmentación de datos causada por la falta de unificación del hardware. \\
Generalización & Seguir completando tareas en escenarios nuevos nunca vistos & Determina si la robótica logra salir del atolladero de la cola larga, como el que vivió la conducción autónoma. \\
Scaling Law & Relación predecible entre datos, tamaño del modelo y desempeño & Si la robótica también presenta una ley de escala, el ritmo de la industria será más claro. \\
\bottomrule
\end{tabularx}
\end{knowledgebox}

\subsection{Por qué los robots siguen atascados en los «escenarios»}

El programa menciona que hoy los escenarios visibles de la robótica siguen concentrados en demostraciones o primeras implementaciones como doblar ropa, preparar café y tareas domésticas. El hardware también gana cada vez más importancia; algunos investigadores incluso creen que podría explicar una proporción muy grande de los factores de éxito. Esto nos recuerda que la robótica no es un problema puramente de modelos. Es a la vez un problema de percepción, control, hardware, costo, confiabilidad, seguridad, cadena de suministro, elección de escenarios y operación posventa.

\begin{warningbox}{Un robot no es «ponerle manos y pies a un LLM»}
La ruta de éxito de los modelos de lenguaje no puede copiarse directamente a la robótica. Los robots deben lidiar con la fricción física, las diferencias de hardware, el ruido de los sensores, los riesgos de seguridad, los escenarios de cola larga y los costos del mundo real. Sin generalización ni aprendizaje continuo, los robots dependerán durante mucho tiempo de la recolección manual de datos y de la ingeniería de escenarios.
\end{warningbox}

\subsection{Resumen de la sección}

Los Neo Labs y Robotics parecen temas dispersos, pero en realidad giran en torno a los datos y a la eficiencia de aprendizaje de la siguiente etapa. Los nuevos laboratorios buscan un nuevo paradigma más allá de Pre-training y RL; las empresas de robótica buscan convertir la manipulación en el mundo real en datos de los que se pueda aprender. El problema que comparten es cómo lograr que un modelo aprenda capacidades generalizables a partir de una cantidad limitada de datos de alta calidad.
\section{Narrativas de China y Estados Unidos y emprendedores chinos: mercado global, ventaja del talento y capital de riesgo}

La última línea principal pasa de la tecnología y el producto a las narrativas de China y Estados Unidos. Guangmi (广密) considera que el capital estadounidense y la desindustrialización de Estados Unidos han movilizado una gran cantidad de dólares, bonos del Tesoro y capital hacia la IA. China, en cambio, se parece más a un rezagado que convierte en industria el 0-1 de Silicon Valley: tiene ventajas en ingeniería, cadena de suministro y talento, pero en los últimos años le han faltado un ambiente de innovación 0-1 y capital de riesgo orientado al mercado. No hace falta tomar este juicio como una gran conclusión; es más útil como marco de decisión para emprendedores.

\begin{figure}[H]
\centering
\includegraphics[width=0.96\textwidth]{china-founder-operating-model.png}
\caption{Modelo de emprendimiento en IA de fundadores chinos: mercado global, talento chino, capital de riesgo, que el producto hable por sí mismo y objetivo de largo plazo. Diagrama conceptual de elaboración propia, a partir de la conversación de 01:10:16--01:18:03.}
\end{figure}

\begin{knowledgebox}{Lectura de la figura: primero el mercado, después la identidad}
La figura empieza por el mercado global, porque las regiones con alta disposición a pagar determinan el techo de ingresos; el talento y la capacidad de ingeniería de China determinan el costo y la velocidad de iteración; el capital de riesgo sostiene el paso de 0 a 1; que el producto hable por sí mismo significa restar peso a las etiquetas de identidad y geopolíticas; y el objetivo de largo plazo es que surjan más empresas tecnológicas multinacionales capaces de tener éxito a la vez en el mercado chino y en el global.
\end{knowledgebox}

\subsection{Diferencias entre las narrativas de IA de China y Estados Unidos}

Esta sección descompone primero la narrativa macro en estructura de recursos, en lugar de quedarse en juicios emocionales. La diferencia entre China y Estados Unidos no es solo «quién es más optimista», sino una combinación distinta de capital, manufactura, mercado, talento y vías de industrialización.

Guangmi describe a Estados Unidos como «impulsado por el capital, impulsado por los datos, impulsado por los Token»: el capital persigue la cúspide y la desindustrialización es grave, pero si la IA tiene éxito, Estados Unidos será muy fuerte. Su descripción de China se parece más a la ventaja del que llega después a la industrialización: muchos prototipos de tecnología y producto pueden surgir en Silicon Valley, pero las empresas chinas pueden alcanzarlos y superarlos apoyándose en ingeniería, cadena de suministro, mercado y capacidad de aplicación. Los vehículos eléctricos, los teléfonos inteligentes, las impresoras 3D y las cámaras deportivas se usan como analogías.

La pregunta implícita en esta analogía es: ¿seguirá la IA el mismo camino? Si el cuello de botella central de la IA es la capacidad de cómputo, los mejores investigadores y el siguiente paradigma, la ventaja de Estados Unidos es más evidente; si el cuello de botella central se desplaza hacia la ingeniería, la aplicación, el costo, los escenarios de uso y la operación global, los equipos chinos podrían recuperar mayores oportunidades.

\subsection{Tres recomendaciones para emprendedores chinos}

Después de discutir las diferencias macro, esta sección pasa a las decisiones que un emprendedor puede ejecutar. Para un equipo individual, lo más importante no es tomar partido en una gran narrativa, sino ordenar correctamente mercado, talento, capital y demostración del producto.

Las recomendaciones del episodio para emprendedores chinos son directas. Primero, apostar con firmeza por el mercado global, en especial por los mercados con alta capacidad de pago. Segundo, aprovechar a los ingenieros y las ventajas industriales de China, sin renunciar a la ventaja del talento. Tercero, restar peso a las etiquetas de identidad y geopolíticas, hacer un buen producto y dejar que el producto hable por sí mismo. Para una empresa en etapa temprana, conviene aceptar todo el dinero que se pueda conseguir, porque la cadena de confianza transfronteriza, los VC estadounidenses de primer nivel y los factores de identidad y geopolíticos aumentan la dificultad.

\begin{importantbox}{El orden realista del emprendedor}
Primero se demuestra el valor del producto; después se busca capital y una mejor narrativa. Presentar una identidad no sustituye al producto, y la narrativa geopolítica tampoco sustituye los ingresos, la retención, la eficiencia ni el pago real de los usuarios.
\end{importantbox}

\subsection{Por qué se espera que China tenga su propio Silicon Valley}

Esta sección sube un nivel desde las recomendaciones para emprendedores y analiza el ecosistema de innovación en sí. Lo que aquí se llama «que China tenga su propio Silicon Valley» no es una copia geográfica, sino el deseo de que los jóvenes, el capital de riesgo, la capacidad de cómputo, la cultura de investigación y el mercado global formen una retroalimentación positiva más densa.

El tono con que cierra el episodio no es una consigna nacionalista, sino preocupación por el ecosistema de innovación. Guangmi recuerda el ambiente de innovación 0-1 de los primeros años del internet móvil en China: el pago móvil, los códigos QR, las transmisiones en vivo, WeChat, el O2O y otros hicieron sentir muy orgullosos a los emprendedores chinos. En los últimos años han surgido muchos unicornios de aplicaciones de IA en el mundo, pero la proporción de equipos chinos es menor que en la era de internet, y eso le parece una lástima.

Considera que la innovación de alto riesgo de los jóvenes necesita el respaldo del capital de riesgo. La base de talento es muy sólida, pero el talento no es el único factor; también se necesitan capital, capacidad de cómputo, un entorno de innovación, tolerancia al fracaso, experiencia en el mercado global y una cultura de investigación de primer nivel. La probabilidad de que dentro de tres a cinco años las empresas de IA más avanzadas sean creadas por equipos chinos es mayor que antes, pero todavía depende de que estas variables puedan completarse.

\subsection{Resumen de la sección}

El valor práctico de las narrativas de China y Estados Unidos consiste en ayudar a los emprendedores a ubicar sus recursos: Estados Unidos tiene capital, investigación de frontera y un mercado con alta capacidad de pago; China tiene ingeniería, talento, cadena de suministro y capacidad de aplicación. La verdadera oportunidad de los emprendedores chinos no está en la narrativa de identidad, sino en lograr combinar el mercado global, el talento chino y la capacidad de ejecución de producto.
\section{Términos clave: índice de palabras clave de este episodio}

Este episodio tiene una densidad de términos muy alta, y muchos de ellos cambian de significado según el contexto. La siguiente tabla reúne los términos centrales para ayudar al lector a repasar.

\begin{longtable}{p{0.23\textwidth}p{0.34\textwidth}p{0.34\textwidth}}
\toprule
Término & Explicación en una frase & Papel en este episodio \\
\midrule
AI War & La competencia en IA como una carrera armamentista que nadie puede permitirse perder & Sustituye el relato de una simple burbuja y explica por qué los gigantes tecnológicos y los Estados siguen invirtiendo. \\
AI Bubble & Debate sobre si las valuaciones y el gasto de capital están sobrecalentados & Pregunta de apertura del programa, pero no el marco de análisis definitivo. \\
SOTA & State of the art, el nivel más avanzado actual & Se advierte que no basta con mirar los benchmarks: también hay que considerar usuarios e ingresos. \\
Benchmark & Tabla de clasificación o métrica que mide la capacidad de un modelo con un conjunto de prueba & Tiene valor de referencia, pero sus puntuaciones pueden inflarse artificialmente o alejarse de las tareas reales. \\
ARR & Annual Recurring Revenue, ingresos recurrentes anuales & Mide la velocidad de comercialización de las empresas AI Native. \\
DAU/MAU & Usuarios activos diarios/mensuales & Compara la frecuencia de uso de ChatGPT y Gemini y su lugar en la mente de los usuarios como punto de entrada. \\
Preentrenamiento & Entrenar un modelo base con cantidades masivas de datos & Se compara con el petróleo: abundante, pero cada vez más cerca de su límite. \\
RL/Post-training & Seguir moldeando el modelo con retroalimentación, preferencias y trayectorias de expertos & Se compara con las energías nuevas: útil, pero limitado en volumen total y en costo. \\
Online Learning & Aprender de la interacción continua y mejorar el comportamiento futuro & Tercer paradigma de este episodio y variable decisiva para cambiar el panorama. \\
Continuous Learning & Aprendizaje continuo; suele usarse como casi sinónimo de Online Learning & Pone énfasis en que el modelo absorba nuevas experiencias a largo plazo. \\
Memory & Capacidad de conservar el contexto del usuario, de las tareas y del historial & Es una de las infraestructuras base de los Agents proactivos y del Online Learning. \\
LoRA & Low-Rank Adaptation, adaptación eficiente del modelo mediante parámetros de rango bajo & El programa la menciona como una de las infraestructuras del nuevo paradigma. \\
MoE & Mixture of Experts, activa solo una parte de las redes expertas según la entrada & Explica que los parámetros activados del modelo no necesariamente crecen sin límite, mientras aumenta la dispersión. \\
Agentic Web & Al convertirse el Agent en punto de entrada, se reorganizan las páginas web, las herramientas, los pagos y la distribución de las plataformas & Afecta la publicidad en buscadores, las Super Apps y las comisiones por transacción. \\
Modelo como producto & La experiencia de producto que percibe el usuario depende cada vez más de la capacidad del modelo & Muestra que una mejora del modelo cambia directamente el valor de las aplicaciones. \\
Datos como modelo & Las diferencias entre modelos provienen, en última instancia, de la distribución de los datos y del ciclo cerrado de retroalimentación & Muestra que las empresas emergentes deben contar con datos diferenciados. \\
Destilación de expertos & Convertir el proceso de trabajo de expertos de alto nivel en datos de entrenamiento & Elemento decisivo para extender horizontalmente las capacidades a distintas industrias. \\
L3/L4 local & Analogía con los niveles de la conducción autónoma para describir una automatización alta en escenarios específicos & Un objetivo comercial más realista que una AGI integral. \\
Neo Labs & Grupo de nuevos laboratorios surgidos de grandes empresas y de OpenAI & Exploran el siguiente paradigma, los sistemas multiagente, los materiales y los modelos de código abierto. \\
Robotics & Inteligencia robótica y sistemas de ejecución física & Dirección del mundo real que pone a prueba la generalización y la coordinación entre datos y hardware. \\
\bottomrule
\end{longtable}

\subsection{Resumen de la sección}

Todos estos términos apuntan a lo mismo: la competencia en IA ya no se limita a la capacidad de un modelo aislado, sino que se ha ampliado a un sistema compuesto de datos, producto, puntos de entrada, capacidad de cómputo, organización, capital y estrategia nacional. Entender cómo se relacionan los términos importa más que memorizar la clasificación de un trimestre.
\section{Síntesis y ampliación}

\subsection{Conclusiones centrales}

Esta última sección condensa todo el texto, que parte de un relato sobre la industria, en criterios de juicio reutilizables. No hace falta recordar las cifras de cada empresa, pero sí conviene quedarse con un mapa por niveles para evaluar la competencia en IA: cómo se restringen mutuamente el gasto de capital, el paradigma de los modelos, los puntos de acceso, los datos y el ecosistema emprendedor.

La línea principal de este episodio puede condensarse en nueve conclusiones.

\begin{enumerate}
\item AI Bubble solo explica el sentimiento detrás de las valuaciones; AI War es lo que explica por qué los gigantes tecnológicos y los países siguen invirtiendo.
\item Los ingresos visibles de OpenAI provienen de las suscripciones, la publicidad y el comercio electrónico, y la API; los ingresos invisibles provienen de los Agent que sustituyen trabajo humano y de la IA que realiza tareas incrementales de alto valor.
\item Las GPU de Nvidia y las TPU de Google representan dos rutas de hardware, nube y modelos: el ecosistema abierto y la integración vertical.
\item Que GPT, Claude y Gemini se alternen el liderazgo se volverá lo normal, hasta que el siguiente paradigma abra de verdad una brecha.
\item Las empresas de modelos fundacionales se parecen a un comercio electrónico generalista; Scale Data es muy eficiente, y las empresas emergentes deben encontrar datos y flujos de trabajo diferenciados.
\item El regreso triunfal de Google no significa que la crisis haya pasado, y ChatGPT tampoco es solo una herramienta de chat; ambos compiten por el punto de acceso de la búsqueda y del asistente personal.
\item El preentrenamiento es como el petróleo, los datos de expertos para RL son como las energías nuevas y el Online Learning es como la fusión nuclear, que aún no logra su avance decisivo.
\item «El modelo es el producto, los datos son el modelo» significa que la ventaja defensiva de un producto de IA debe volver a la distribución de los datos, la retroalimentación de los usuarios y los procesos de los expertos.
\item Las oportunidades de los emprendedores de IA de origen chino provienen del mercado global, del talento chino, de la capacidad de ingeniería y de dejar que el producto hable por sí mismo, no de presentarse a partir de su identidad.
\end{enumerate}

\subsection{Preguntas abiertas}

Estas preguntas determinarán la bifurcación a partir de 2026.

\begin{itemize}
\item ¿El Online Learning dará señales claras en 2026, o seguirá siendo un relato de investigación?
\item ¿ChatGPT integrará la búsqueda más rápido de lo que Google logra renovarse a sí mismo?
\item ¿Podrán la multimodalidad de Gemini y el ecosistema Android convertirse en un lugar en la mente del usuario como asistente personal de uso frecuente?
\item ¿Podrán el Coding y la Beta de la estrategia empresarial de Anthropic seguir resistiendo la expansión horizontal de OpenAI y Google?
\item En el Agentic Web, ¿quién estandarizará primero los pagos, la autorización, la verificación de humano frente a máquina y los protocolos de los sitios web?
\item ¿Convergerán las rutas de datos de la Robotics, o seguirán divergiendo a largo plazo?
\item ¿Podrán los equipos chinos recuperar, entre los unicornios globales de aplicaciones de IA, una proporción como la que tuvieron en la era de internet?
\end{itemize}

\subsection{Lecturas adicionales}

\begin{itemize}
\item EP136, informe trimestral de los grandes modelos a nivel global, episodio 9: Coding, el modelo como OS y los tres grandes de Silicon Valley.
\item EP134, panorama de los datos: la pirámide de datos, Recipe, el etiquetado y la fijación de precios.
\item EP139, panorama de la historia técnica de los Agent: la genealogía técnica de los Agent y su alcance social.
\item EP130, entrevista con Zhang Yueguang (张月光): productos AI Native, diseño del contexto y metodología de producto para Agent.
\item Entrevistas públicas e informes técnicos relacionados con OpenAI, Anthropic, Google DeepMind, SSI y Thinking Machines: para entender el siguiente paradigma y la ruta del Post-training.
\end{itemize}

\begin{importantbox}{El juicio final}
Lo que más vale la pena conservar de EP127 es un marco que atraviesa niveles: la competencia en IA no es solo una tabla de clasificación de modelos ni solo una burbuja de capital, sino un cambio simultáneo en el cómputo, los datos, los puntos de acceso, los productos, los paradigmas de aprendizaje y las estrategias nacionales. El ganador futuro no será necesariamente la empresa con el mejor benchmark de algún trimestre, sino el sistema capaz de cerrar el ciclo que une datos reales, aprendizaje continuo y puntos de acceso de los usuarios.
\end{importantbox}

\end{document}
